\documentclass[10pt,letterpaper]{article}
\usepackage{cornell-notes}

\title{Clause 28.07.03: Three Dimensional Hypotenuse}
\author{}
\date{}
\setDocTitle{Clause 28.07.03: Three Dimensional Hypotenuse}
\setDocSubtitle{C++ 2024}
\setDocOwner{C++ 2024 Cornell Notes}
\setDocDate{\today}
\setCornellCollection{C++ 2024 Cornell Notes}
\setCornellUnitType{Clause}
\setCornellUnitNumber{28.07.03}
\setCornellUnitTitle{Three Dimensional Hypotenuse}
\hypersetup{pdftitle={Clause 28.07.03: Three Dimensional Hypotenuse},pdfsubject={Cornell notes},pdfauthor={}}

\begin{document}
\maketitle
\begin{CornellOverviewBox}{Overview}
\textbf{Classification:} Standard library - Normative\\[2pt]
\textbf{Learning objective:} Understand the rules, terminology, and semantic consequences associated with Three-dimensional hypotenuse.
\end{CornellOverviewBox}\begin{CornellSummaryBox}{Summary}
{\sffamily\bfseries\color{Accent} Summary}\par\vspace{3pt}Section 28.7.3 focuses on Three-dimensional hypotenuse. Apply the specified interface together with its mandates, effects, result conditions, exception behavior, and complexity constraints. When applying it, preserve the distinction between mandatory requirements, recommendations, permissions, and behavior deliberately left undefined, unspecified, or implementation-defined.
\end{CornellSummaryBox}

\end{document}
